\section{Introduction}
Large Language Models (LLMs) have attracted ever-growing attention in the past few years due to their outstanding ability in answering complex questions and generating human-quality text~\cite{vaswani2017attention,wang2024current}.
Despite remarkable progress, LLMs still exhibit significant limitations when faced with complex tasks that require integrating vast and diverse knowledge domains or require deep reasoning and multi-step logical inference~\cite{talmor2019commonsenseqa,petroni2021kilt}. 
As shown in Figure~\ref{fig:example}(a), LLMs tend to provide an incorrect response when the question
requires knowledge beyond their scope.  
Graph Retrieval-Augmented Generation (GraphRAG)~\cite{xiang2025use,zhangsurvey,pengsurvey,hansurvey,zhu2024structugraphrag,li2025g}, which utilizes knowledge graphs for knowledge representation, enables effective modeling of complex real-world relationships, such as causality, hierarchy, and dependency~\cite{micro_graphrag,pengsurvey}.

However, existing GraphRAG methods still face significant performance bottlenecks when handling complex problems that require multi-hop reasoning. The main bottleneck is that existing methods rely on pre-defined or heuristic rules during query processing and retrieval, which are too rigid and inflexible. Concretely, during query processing, the semantic or lexical similarity search~\cite{hansurvey,gao2023retrieval,fan2024survey} remains the most widely adopted. For semantic similarity search, queries with high complexity are averaged out during vectorization, making it difficult to gather accurate information for each distinct semantic element. On the other hand, lexical similarity search captures only the superficial meaning of the query, easily losing implicit relationships and reasoning logic. Therefore, as shown in Figure~\ref{fig:example}(b), existing GraphRAG methods are unable to answer the query, since the retrieval results are incomplete or incorrect when queries are overly complex.
Besides, during the graph retrieval phase, current methods employ a combination of graph traversal and subgraph retrieval for path reasoning, which are rather rigid~\cite{tog,tog2,jimenez2024hipporag,guo2024lightrag,li2024graphreader}. For instance, ToG~\cite{tog,tog2}, a representative GraphRAG method, starts from a node in the graph and traverses the surrounding nodes based on breadth-first search strategy with pruning. While this strategy might work for short paths, it struggles with long-path problems, which require capturing logically sound but distant or even disconnected relationships. 
This inflexibility often leads to broken reasoning chains and failed retrievals, making it unsuitable for multi-hop problems with complex reasoning logic.
As a result, the existing GraphRAG methods fail to fully leverage the LLM's deep reasoning potential.%


In this paper, we try to further enhance the multi-hop reasoning ability of GraphRAG.
Reinforcement learning (RL) has recently been shown to be effective in endowing LLMs with powerful reasoning abilities. For example, in the case of DeepSeek-R1~\cite{guo2025deepseek}, it has been discovered that large-scale reinforcement learning can endow models with deep thinking capabilities~\cite{xie2025logic,jaech2024openai,team2025kimi}, even without any supervised fine-tuning. This effectiveness may result from RL's use of reward signals to encourage exploration, refining dynamic reasoning strategies from static knowledge~\cite{ouyang2022training} beyond static reasoning patterns~\cite{asai2024self,zhang2025corag,xiong2025mcts,deeprag}. However, how to utilize RL for GraphRAG training remains unexplored in the literature. The major technical obstacle lies in how to design proper reward functions so that the reasoning process of LLMs within GraphRAG can be properly constrained, avoiding potential problems such as shallow retrieval~\cite{helwe2021reasoning,deng2025boosting} and over-thinking~\cite{sui2025stop,chen2024not,cuesta2025large}. 

To address this issue, we propose \methodname, an adaptive GraphRAG framework with deep reasoning capabilities through process-constrained outcome-based RL. As shown in Figure~\ref{fig:example}(c), our method can decompose complex multi-hop problems by reasoning and autonomous invocation of retrieval tools. Specifically, we employ a modified version of Group Relative Policy Optimization (GRPO) that supports rollout-with-thinking mechanism~\cite{song2025r1searcher} to leverage external retrieval tools on demand. Next, we design two process-constrained reward functions, Progressive Retrieval Attenuation (PRA) and Cost-Aware F1 (CAF), aiming to inject prior knowledge into the LLM's training process.
In detail, PRA aims to dynamically adjust the reward for retrieval actions. It encourages necessary retrievals
, preventing the model from shallow retrieval. 
As the reasoning deepens, PRA progressively attenuates retrieval rewards, 
guiding the model to 
avoid unnecessary retrieval overhead. 
On the other hand, CAF evaluates the final answer correctness using the F1 score while also accounting for the computational resources consumed during the reasoning and retrieval process. It encourages the model to achieve correct results more efficiently to prevent over-thinking. These two rewards collectively balance effectiveness and computational cost. We further design a three-stage phase-dependent training strategy. In the first stage, we focus on the answer format to alleviate cold start problems. In the second and third stages, we focus on improving the answer quality using the aforementioned two rewards. 
Lastly, our method employs a hybrid graph-textual retrieval.
By leveraging graph structures to precisely capture entity relationships while preserving the rich semantic details of the original text, our method provides more comprehensive and integrated knowledge for complex reasoning.



Experiments show that \methodname~significantly outperforms existing GraphRAG methods in solving complex reasoning problems, across both in-domain and out-of-domain datasets. Additionally, our framework offers flexible integration with various retrieval methods, consistently enhancing their performance.

Our main contributions are summarized as follows:
\begin{itemize}[leftmargin=0.5cm]
\item We propose \methodname, a pioneering framework that utilizes GRPO that supports rollout-with-thinking to enhance the reasoning ability of GraphRAG. %
\item We propose PRA and CAF, two tailored reward designs to guide the training process of RL and handle the shallow retrieval and over-thinking problems. We also propose a phase-dependent training strategy and a hybrid graph-textual retrieval. 

\item Experiments demonstrate that our method exhibits state-of-the-art performance on both in-domain and out-of-domain datasets, and can be flexibly integrated with various existing retrieval methods with consistently improved performance.
\end{itemize}


\begin{figure*}
  
  \includegraphics[width=\textwidth,page=1]{Figures/example.pdf}
  \vspace{-0.8cm}
  \caption{An example of comparing using only LLMs, GraphRAG, and \methodname~in answering complex problems: (a) The LLM (Qwen2.5-7B) directly produces an output, but the answer is incorrect, (b) HippoRAG2~\cite{gutierrez2025hipporag2} enhances the LLM (Qwen2.5-7B) by retrieving external knowledge, but fail to generate correct outputs for complex problem,  (c) our method successfully decomposes the problem and autonomously invokes retrieval methods, producing the correct output.}
  \Description{This figure provides a comparative schematic illustrating the performance of three distinct approaches in handling the same complex reasoning problem, visually demonstrating the superiority of the proposed GraphRAG-R1 framework. The schematic is divided into three juxtaposed sub-figures:
(a) Using only the base Large Language Model (LLM, Qwen2.5-7B): The model directly generates an answer, which is incorrect. This reveals the limitations of LLMs regarding knowledge scope and multi-step reasoning.
(b) Using an existing Graph Retrieval-Augmented Generation method (HippoRAG2) to augment the same LLM: While this method introduces external knowledge retrieval, it still fails to produce the correct answer for this complex problem. This highlights the bottlenecks of traditional GraphRAG methods in flexible problem decomposition and multi-hop reasoning.
(c) Applying the proposed GraphRAG-R1 method: The model successfully decomposes the complex problem into sub-steps, autonomously and on-demand invokes external retrieval tools to acquire necessary information, and finally integrates the reasoning to produce the correct answer. Through this comparison, the figure clearly validates the effectiveness of GraphRAG-R1 in enhancing the LLM's capability to solve complex problems.}
  \label{fig:example}
\vspace{-0.3cm}
\end{figure*}
